% Lösungen Einheit 1 von 4 – Einheitskreis und Bogenmaß (zusammenbau v0.1)
\einheitenkopf{Einheit 1 von 4 · Einheitskreis und Bogenmaß}

%% TODO Lösung der Erklärzeile 14g) fehlt in der Bank
\erg{14}{a) Bogenmaß und RAD – Dezimalzahl ohne Gradzeichen \quad b) P auf dem Kreis, 40° gegen den Uhrzeigersinn ab der positiven Rechtsachse; P ≈ (0,77|0,64) \quad c) P auf dem Kreis, 130° gegen den Uhrzeigersinn ab der positiven Rechtsachse; P ≈ (−0,64|0,77) \quad d) P auf dem Kreis, 215° gegen den Uhrzeigersinn ab der positiven Rechtsachse; P ≈ (−0,82|−0,57) \quad e) P auf dem Kreis, 320° gegen den Uhrzeigersinn ab der positiven Rechtsachse; P ≈ (0,77|−0,64) \quad f) P auf dem Kreis, 70° gegen den Uhrzeigersinn ab der positiven Rechtsachse; P ≈ (0,34|0,94) \quad g) \ldots \quad h) $\alpha \approx 100$° \quad i) sin $\alpha \approx 0{,}8$; cos $\alpha \approx 0{,}6$ \quad j) $\alpha = 90$°, sin $\alpha = 1$, cos $\alpha = 0$ \quad k) sin 50° ≈ 0,77 – die Ablesung passt auf eine Stelle \quad l) sin positiv, cos negativ – zweiter Quadrant (links oben) \quad m) $\alpha_1 \approx 36{,}9$°; $\alpha_2 \approx 143{,}1$° \quad n) Umfang $2\pi$, ein Sechstel davon: $\frac{2\pi}{6} = \frac{\pi}{3} \approx 1{,}05$ \quad o) $\frac{2}{9}\pi \approx 0{,}70$ \quad p) 135° (drei Viertel von 180°) \quad q) $\frac{\pi}{3}$ \quad r) DEG; im Modus RAD: sin 1 ≈ 0,84 \quad s) e ≈ 8,11 cm; sin 128° ≈ 0,79 ist die Höhe des Punktes zu 128° auf dem Einheitskreis – er liegt links oben, die Höhe ist positiv \quad t) a) ≈ 1,83; b) = 240°; c) ≈ 0,97 – sin 105° ist die Höhe des Punktes zu 105° auf dem Einheitskreis}
\erg{15}{a) Fehler: Der Rechner steht auf RAD, 25 wurde als Bogenmaß gerechnet. Richtig im Modus DEG: sin 25° ≈ 0,42 \quad b) Der Sinuswert ist die Höhe eines Punktes auf dem Kreis mit dem Radius 1 – höher als der Radius kann kein Punkt des Kreises liegen.}
